\documentclass[aspectratio=169]{beamer}
\usetheme{metropolis}
\usepackage{graphicx}
\usepackage{tikz}
\usepackage{amsmath}
\usepackage{listings}
\usepackage{xcolor}
\usepackage{booktabs}
\usepackage{emoji}

% Code listing setup
\lstset{
  basicstyle=\ttfamily\small,
  keywordstyle=\color{blue},
  commentstyle=\color{gray},
  stringstyle=\color{red},
  showstringspaces=false,
  breaklines=true,
  frame=single,
  language=Python
}

\title{Dimensionality Reduction \& Text Embeddings}
\subtitle{Weeks 3-4: From Distributional Semantics to Word Vectors}
\author{LLM Course}
\date{\today}

\begin{document}

% ============================================================
% TITLE SLIDE
% ============================================================
\begin{frame}
\titlepage
\end{frame}

% ============================================================
% OVERVIEW
% ============================================================
\begin{frame}{Overview 📋}
\large
\begin{enumerate}
    \item 🧠 \textbf{Distributional Semantics}: The Foundation
    \item 📊 \textbf{Dimensionality Reduction}: PCA, UMAP, t-SNE
    \item 📚 \textbf{Classic Methods}: LSA \& LDA
    \item 🎨 \textbf{Modern Topic Modeling}: Top2Vec, BERTopic
    \item 🚀 \textbf{Neural Word Embeddings}: Word2Vec, GloVe, FastText
    \item 🔮 \textbf{Preview}: Transformer Embeddings
    \item ✍️ \textbf{Assignment 3}: Hands-on Practice
\end{enumerate}

\vspace{1em}
\textit{Goal: Understand how machines learn to represent meaning}
\end{frame}

% ============================================================
% FUNDAMENTAL QUESTION
% ============================================================
\begin{frame}{The Fundamental Question 🤔}
\begin{center}
\LARGE
\textbf{How do we represent the meaning of words in a way that computers can understand?}
\end{center}

\pause
\vspace{2em}

\begin{columns}[T]
\begin{column}{0.48\textwidth}
\textbf{Traditional Approach:}
\begin{itemize}
    \item Dictionaries
    \item Taxonomies (WordNet)
    \item Manual feature engineering
\end{itemize}
\end{column}

\begin{column}{0.48\textwidth}
\textbf{Distributional Approach:}
\begin{itemize}
    \item Learn from data
    \item Context-based
    \item Scalable
\end{itemize}
\end{column}
\end{columns}
\end{frame}

% ============================================================
% DISTRIBUTIONAL HYPOTHESIS
% ============================================================
\begin{frame}{The Distributional Hypothesis 📖}
\begin{center}
\huge
\textit{"You shall know a word by the company it keeps"}
\end{center}

\begin{center}
--- J.R. Firth (1957)
\end{center}

\pause
\vspace{2em}

\textbf{Key Idea:}
Words that appear in similar contexts tend to have similar meanings.

\vspace{1em}

\begin{exampleblock}{Example}
\begin{itemize}
    \item "The \textbf{cat} sat on the mat"
    \item "The \textbf{dog} sat on the mat"
    \item "The \textbf{kitten} sat on the mat"
\end{itemize}
$\rightarrow$ \textit{cat}, \textit{dog}, \textit{kitten} are semantically related
\end{exampleblock}

\vspace{0.5em}
\small
\textit{References: Boleda (2020), Grand et al. (2022)}
\end{frame}

% ============================================================
% COGNITIVE VS COMPUTATIONAL
% ============================================================
\begin{frame}{Cognitive vs. Computational Perspectives 🧠💻}
\begin{columns}[T]
\begin{column}{0.48\textwidth}
\textbf{Human Cognition:}
\begin{itemize}
    \item Grounded in experience
    \item Multimodal (vision, sound, touch)
    \item Embodied understanding
    \item Cultural context
    \item Emotional associations
\end{itemize}

\vspace{1em}
\alert{How do YOU represent "coffee"?}
\end{column}

\pause

\begin{column}{0.48\textwidth}
\textbf{Computational Models:}
\begin{itemize}
    \item Patterns in text data
    \item Co-occurrence statistics
    \item Vector representations
    \item Distributional semantics
    \item Learned from corpora
\end{itemize}

\vspace{1em}
\alert{Model represents "coffee" by its linguistic context}
\end{column}
\end{columns}

\vspace{1em}
\begin{center}
\textbf{Question:} What does \textit{semantic similarity} really mean?
\end{center}

\small
\textit{Reference: Grand et al. (2022) - "Semantic projection recovers rich human knowledge"}
\end{frame}

% ============================================================
% VECTOR SPACE MODEL
% ============================================================
\begin{frame}{From Words to Vectors 🗺️}
\textbf{Goal:} Represent each word as a point in high-dimensional space

\vspace{1em}

\begin{center}
\begin{tikzpicture}[scale=0.8]
    % Axes
    \draw[->] (0,0) -- (6,0) node[right] {Context: "furry"};
    \draw[->] (0,0) -- (0,5) node[above] {Context: "barks"};

    % Points
    \node[circle,fill=blue,inner sep=2pt,label=right:{dog}] at (4.5,4) {};
    \node[circle,fill=blue,inner sep=2pt,label=right:{puppy}] at (4,3.5) {};
    \node[circle,fill=red,inner sep=2pt,label=left:{cat}] at (3.5,1) {};
    \node[circle,fill=red,inner sep=2pt,label=left:{kitten}] at (3,0.8) {};
    \node[circle,fill=green,inner sep=2pt,label=below:{car}] at (0.5,0.5) {};
    \node[circle,fill=green,inner sep=2pt,label=below:{vehicle}] at (0.7,0.3) {};

    % Grid
    \draw[help lines, gray!30] (0,0) grid (6,5);
\end{tikzpicture}
\end{center}

\textbf{Problem:} Real vocabularies have 10,000+ words, contexts are even more numerous!

$\rightarrow$ We need \textbf{dimensionality reduction} 📉
\end{frame}

% ============================================================
% DIMENSIONALITY REDUCTION INTRO
% ============================================================
\begin{frame}{Why Dimensionality Reduction? 📊}
\begin{columns}[T]
\begin{column}{0.48\textwidth}
\textbf{The Curse of Dimensionality:}
\begin{itemize}
    \item Vocabulary: 50,000 words
    \item Context features: 50,000+
    \item Matrix: 50k × 50k = 2.5B entries!
    \item Sparse (mostly zeros)
    \item Computationally expensive
    \item Hard to visualize
\end{itemize}
\end{column}

\begin{column}{0.48\textwidth}
\textbf{Solution:}
\begin{itemize}
    \item Compress to \textasciitilde50-300 dimensions
    \item Preserve semantic relationships
    \item Dense representations
    \item Computationally efficient
    \item Easier to work with
    \item Can visualize in 2D/3D!
\end{itemize}
\end{column}
\end{columns}

\vspace{1em}
\begin{center}
\Large
50,000 dimensions $\rightarrow$ 100 dimensions ✨
\end{center}
\end{frame}

% ============================================================
% PCA
% ============================================================
\begin{frame}{Principal Component Analysis (PCA) 📈}
\textbf{Intuition:} Find the directions of maximum variance

\vspace{1em}

\begin{columns}[T]
\begin{column}{0.48\textwidth}
\begin{center}
\begin{tikzpicture}[scale=0.9]
    % Data points
    \foreach \x/\y in {1/1, 1.5/1.3, 2/1.8, 2.5/2.2, 3/2.5, 3.5/3, 4/3.3, 4.5/3.8}
        \fill[blue] (\x,\y) circle (2pt);

    % Principal component
    \draw[->, very thick, red] (0.5,0.5) -- (5,4.2) node[above right] {PC1};
    \draw[->, thick, orange] (2.5,2) -- (1.5,3.5) node[above] {PC2};

    % Axes
    \draw[->] (0,0) -- (5.5,0) node[right] {$x_1$};
    \draw[->] (0,0) -- (0,5) node[above] {$x_2$};
\end{tikzpicture}
\end{center}
\end{column}

\begin{column}{0.48\textwidth}
\textbf{Algorithm:}
\begin{enumerate}
    \item Center the data
    \item Compute covariance matrix
    \item Find eigenvectors
    \item Project onto top k eigenvectors
\end{enumerate}

\vspace{0.5em}
\textbf{Properties:}
\begin{itemize}
    \item Linear transformation
    \item Preserves global structure
    \item Fast \& interpretable
\end{itemize}
\end{column}
\end{columns}

\vspace{1em}
\textbf{Limitation:} Assumes linear relationships (often not true for text!)
\end{frame}

% ============================================================
% UMAP
% ============================================================
\begin{frame}{UMAP: Uniform Manifold Approximation \& Projection 🗺️}
\textbf{Better for visualizing complex, non-linear structure in text data}

\vspace{0.5em}

\begin{columns}[T]
\begin{column}{0.48\textwidth}
\textbf{Key Advantages:}
\begin{itemize}
    \item Preserves both local \& global structure
    \item Non-linear dimensionality reduction
    \item Faster than t-SNE
    \item Theoretically grounded
    \item Great for visualization
\end{itemize}

\vspace{0.5em}
\textbf{Use Cases:}
\begin{itemize}
    \item Visualizing word embeddings
    \item Topic modeling
    \item Clustering analysis
    \item Exploratory data analysis
\end{itemize}
\end{column}

\begin{column}{0.48\textwidth}
\begin{lstlisting}[language=Python]
import umap
from sklearn.datasets import load_digits

# Example: Reduce to 2D
reducer = umap.UMAP(
    n_components=2,
    n_neighbors=15,
    min_dist=0.1
)

embedding = reducer.fit_transform(
    word_vectors
)

# Plot
plt.scatter(
    embedding[:, 0],
    embedding[:, 1],
    c=labels
)
\end{lstlisting}
\end{column}
\end{columns}

\vspace{0.5em}
\small
\textit{Reference: McInnes \& Healy (2018) - UMAP: Uniform Manifold Approximation and Projection}
\end{frame}

% ============================================================
% PCA VS UMAP
% ============================================================
\begin{frame}{PCA vs. UMAP vs. t-SNE 🔍}
\begin{center}
\begin{tabular}{lccc}
\toprule
\textbf{Property} & \textbf{PCA} & \textbf{t-SNE} & \textbf{UMAP} \\
\midrule
Linear/Non-linear & Linear & Non-linear & Non-linear \\
Speed & Fast ⚡ & Slow 🐌 & Fast ⚡ \\
Global structure & ✓ & ✗ & ✓ \\
Local structure & $\sim$ & ✓ & ✓ \\
Deterministic & ✓ & ✗ & $\sim$ \\
Scalability & Excellent & Poor & Good \\
\bottomrule
\end{tabular}
\end{center}

\vspace{1em}

\begin{block}{Recommendation}
\begin{itemize}
    \item \textbf{PCA}: First pass, quick exploration, preprocessing
    \item \textbf{t-SNE}: 2D/3D visualization (small datasets)
    \item \textbf{UMAP}: Best of both worlds - use for most NLP tasks!
\end{itemize}
\end{block}
\end{frame}

% ============================================================
% LSA
% ============================================================
\begin{frame}{Latent Semantic Analysis (LSA) 📚}
\textbf{The OG of semantic embeddings (1990)}

\vspace{1em}

\begin{columns}[T]
\begin{column}{0.48\textwidth}
\textbf{Algorithm:}
\begin{enumerate}
    \item Build term-document matrix $X$
    \item Apply TF-IDF weighting
    \item Perform SVD: $X = U\Sigma V^T$
    \item Keep top $k$ dimensions
    \item Use $U_k$ as word embeddings
\end{enumerate}

\vspace{0.5em}
\textbf{Key Innovation:}
\begin{itemize}
    \item Discovers latent topics
    \item Solves synonymy problem
    \item Matrix factorization
\end{itemize}
\end{column}

\begin{column}{0.48\textwidth}
\begin{center}
\begin{tikzpicture}[scale=0.7]
    % Term-Document Matrix
    \draw (0,0) rectangle (2,3);
    \node at (1, 3.5) {Terms};
    \node[rotate=90] at (-0.5, 1.5) {Docs};
    \node at (1, 1.5) {$X$};

    % SVD
    \node at (3, 1.5) {$=$};

    % U
    \draw (4,0) rectangle (5,3);
    \node at (4.5, 1.5) {$U$};

    % Sigma
    \draw (5.5,1.5) rectangle (7.5,3);
    \node at (6.5, 2.25) {$\Sigma$};

    % V^T
    \draw (8,1.5) rectangle (10,3);
    \node at (9, 2.25) {$V^T$};
\end{tikzpicture}
\end{center}

\vspace{0.5em}
\textbf{Limitations:}
\begin{itemize}
    \item Linear relationships only
    \item Bag-of-words (no order)
    \item Computationally expensive
    \item No polysemy handling
\end{itemize}
\end{column}
\end{columns}

\vspace{0.5em}
\small
\textit{Reference: Deerwester et al. (1990) - Indexing by Latent Semantic Analysis}
\end{frame}

% ============================================================
% LDA
% ============================================================
\begin{frame}{Latent Dirichlet Allocation (LDA) 🎲}
\textbf{Probabilistic topic modeling}

\vspace{0.5em}

\begin{columns}[T]
\begin{column}{0.48\textwidth}
\textbf{Generative Story:}
\begin{enumerate}
    \item For each document:
    \begin{itemize}
        \item Draw topic distribution $\theta \sim \text{Dir}(\alpha)$
    \end{itemize}
    \item For each word in document:
    \begin{itemize}
        \item Choose topic $z \sim \text{Mult}(\theta)$
        \item Choose word $w \sim \text{Mult}(\phi_z)$
    \end{itemize}
\end{enumerate}

\vspace{0.5em}
\textbf{Key Ideas:}
\begin{itemize}
    \item Documents are mixtures of topics
    \item Topics are distributions over words
    \item Bayesian inference
\end{itemize}
\end{column}

\begin{column}{0.48\textwidth}
\begin{lstlisting}[language=Python]
from sklearn.decomposition import LatentDirichletAllocation

# Fit LDA model
lda = LatentDirichletAllocation(
    n_components=10,  # topics
    random_state=42
)

doc_topics = lda.fit_transform(
    doc_term_matrix
)

# Print top words per topic
for idx, topic in enumerate(
    lda.components_
):
    top_words = [
        vocab[i] for i in
        topic.argsort()[-10:]
    ]
    print(f"Topic {idx}: {top_words}")
\end{lstlisting}
\end{column}
\end{columns}

\vspace{0.5em}
\small
\textit{Reference: Blei et al. (2003) - Latent Dirichlet Allocation}
\end{frame}

% ============================================================
% LDA EXAMPLE
% ============================================================
\begin{frame}{LDA Example: Topic Discovery 🔍}
\begin{center}
\textbf{Example Topics from News Corpus:}
\end{center}

\vspace{1em}

\begin{block}{Topic 1: Politics}
\texttt{government, election, president, vote, policy, senate, congress, party}
\end{block}

\begin{block}{Topic 2: Sports}
\texttt{game, team, player, score, win, season, coach, championship}
\end{block}

\begin{block}{Topic 3: Technology}
\texttt{software, computer, data, algorithm, system, internet, code, AI}
\end{block}

\begin{block}{Topic 4: Finance}
\texttt{market, stock, price, investment, economy, trade, bank, profit}
\end{block}

\vspace{0.5em}
\textbf{Document representation:} [0.7, 0.1, 0.15, 0.05] $\leftarrow$ mostly politics
\end{frame}

% ============================================================
% MODERN TOPIC MODELING
% ============================================================
\begin{frame}{Modern Topic Modeling 🎨}
\textbf{Combining neural embeddings with topic models}

\vspace{1em}

\begin{columns}[T]
\begin{column}{0.48\textwidth}
\textbf{Top2Vec:}
\begin{itemize}
    \item Embed documents with Doc2Vec
    \item Cluster in embedding space
    \item Topics = cluster centroids
    \item Automatic topic number detection
    \item No stopword removal needed!
\end{itemize}

\vspace{0.5em}
\begin{lstlisting}[language=Python]
from top2vec import Top2Vec

model = Top2Vec(
    documents,
    embedding_model='doc2vec'
)

topics = model.get_topics()
\end{lstlisting}
\end{column}

\begin{column}{0.48\textwidth}
\textbf{BERTopic:}
\begin{itemize}
    \item Embed with transformers (BERT)
    \item Reduce with UMAP
    \item Cluster with HDBSCAN
    \item Topics via TF-IDF on clusters
    \item State-of-the-art results!
\end{itemize}

\vspace{0.5em}
\begin{lstlisting}[language=Python]
from bertopic import BERTopic

model = BERTopic(
    language="english",
    calculate_probabilities=True
)

topics, probs = model.fit_transform(
    documents
)
\end{lstlisting}
\end{column}
\end{columns}

\vspace{0.5em}
\textbf{Advantage:} Leverage power of pre-trained neural networks!
\end{frame}

% ============================================================
% WORD2VEC INTRO
% ============================================================
\begin{frame}{Word2Vec: The Revolution 🚀}
\begin{center}
\LARGE
\textbf{2013: Everything changed}
\end{center}

\vspace{1em}

\begin{columns}[T]
\begin{column}{0.48\textwidth}
\textbf{Key Innovation:}
\begin{itemize}
    \item Neural network approach
    \item Predict context from word (CBOW)
    \item Predict word from context (Skip-gram)
    \item Dense, low-dimensional vectors
    \item Captures semantic relationships
    \item Fast training
\end{itemize}
\end{column}

\begin{column}{0.48\textwidth}
\textbf{Famous Result:}
\begin{center}
\Large
$\vec{king} - \vec{man} + \vec{woman} \approx \vec{queen}$
\end{center}

\vspace{1em}
\textbf{Other Examples:}
\begin{itemize}
    \item $\vec{Paris} - \vec{France} + \vec{Italy} \approx \vec{Rome}$
    \item $\vec{walking} - \vec{walk} + \vec{swim} \approx \vec{swimming}$
\end{itemize}
\end{column}
\end{columns}

\vspace{1em}
\small
\textit{Reference: Mikolov et al. (2013a,b) - Efficient Estimation of Word Representations}
\end{frame}

% ============================================================
% WORD2VEC ARCHITECTURES
% ============================================================
\begin{frame}{Word2Vec Architectures 🏗️}
\begin{columns}[T]
\begin{column}{0.48\textwidth}
\textbf{CBOW (Continuous Bag of Words):}

\vspace{0.5em}
\begin{center}
\begin{tikzpicture}[scale=0.7]
    % Input layer
    \node at (-2, 4) {the};
    \node at (-1, 4) {cat};
    \node at (1, 4) {on};
    \node at (2, 4) {the};

    % Arrows down
    \foreach \x in {-2,-1,1,2}
        \draw[->] (\x, 3.7) -- (0, 2.3);

    % Hidden layer
    \draw[fill=blue!20] (-0.5,1.5) rectangle (0.5,2.5);
    \node at (0, 2) {Hidden};

    % Output layer
    \draw[->] (0, 1.5) -- (0, 0.3);
    \draw[fill=green!20] (-0.5,0) rectangle (0.5,0.3);
    \node at (0, -0.4) {\textbf{sat}};
\end{tikzpicture}
\end{center}

Predict center word from context
\end{column}

\begin{column}{0.48\textwidth}
\textbf{Skip-gram:}

\vspace{0.5em}
\begin{center}
\begin{tikzpicture}[scale=0.7]
    % Input layer
    \node at (0, 4) {\textbf{sat}};

    % Arrow down
    \draw[->] (0, 3.7) -- (0, 2.5);

    % Hidden layer
    \draw[fill=blue!20] (-0.5,1.5) rectangle (0.5,2.5);
    \node at (0, 2) {Hidden};

    % Output layer
    \foreach \x in {-2,-1,1,2}
        \draw[->] (0, 1.5) -- (\x, 0.3);

    \node at (-2, 0) {the};
    \node at (-1, 0) {cat};
    \node at (1, 0) {on};
    \node at (2, 0) {the};
\end{tikzpicture}
\end{center}

Predict context from center word
\end{column}
\end{columns}

\vspace{1em}

\textbf{Training Trick:} Negative sampling - don't update all weights!
\begin{itemize}
    \item Sample a few "negative" examples (words that don't appear in context)
    \item Much faster than full softmax over entire vocabulary
\end{itemize}
\end{frame}

% ============================================================
% WORD2VEC CODE
% ============================================================
\begin{frame}[fragile]{Word2Vec in Practice 💻}
\begin{lstlisting}[language=Python]
from gensim.models import Word2Vec
import gensim.downloader as api

# Load pre-trained model
model = api.load('word2vec-google-news-300')

# Find similar words
similar = model.most_similar('computer', topn=5)
# Output: [('computers', 0.72), ('laptop', 0.69), ...]

# Analogy: king - man + woman = ?
result = model.most_similar(
    positive=['king', 'woman'],
    negative=['man'],
    topn=1
)
# Output: [('queen', 0.71)]

# Train your own model
sentences = [['the', 'cat', 'sat'], ['the', 'dog', 'ran']]
custom_model = Word2Vec(
    sentences,
    vector_size=100,
    window=5,
    min_count=1,
    sg=1  # Skip-gram
)
\end{lstlisting}
\end{frame}

% ============================================================
% GLOVE
% ============================================================
\begin{frame}{GloVe: Global Vectors 🌍}
\textbf{Combining count-based \& prediction-based methods}

\vspace{1em}

\begin{columns}[T]
\begin{column}{0.48\textwidth}
\textbf{Key Insight:}
\begin{itemize}
    \item Word2Vec uses local context windows
    \item But global statistics matter too!
    \item Ratios of co-occurrence probabilities are meaningful
\end{itemize}

\vspace{0.5em}
\textbf{Example:}
\begin{itemize}
    \item $P(\text{solid}|\text{ice})$ is high
    \item $P(\text{solid}|\text{steam})$ is low
    \item $\frac{P(\text{solid}|\text{ice})}{P(\text{solid}|\text{steam})}$ is very high!
\end{itemize}
\end{column}

\begin{column}{0.48\textwidth}
\textbf{Objective Function:}
\begin{align*}
J = \sum_{i,j=1}^{V} f(X_{ij}) (\vec{w}_i^T \tilde{\vec{w}}_j + b_i + \tilde{b}_j - \log X_{ij})^2
\end{align*}

where:
\begin{itemize}
    \item $X_{ij}$ = co-occurrence count
    \item $f$ = weighting function
    \item $\vec{w}_i$ = word vector
    \item $\tilde{\vec{w}}_j$ = context vector
\end{itemize}

\vspace{0.5em}
\textbf{Result:}
\begin{itemize}
    \item Similar performance to Word2Vec
    \item More interpretable
    \item Faster training on large corpora
\end{itemize}
\end{column}
\end{columns}

\vspace{0.5em}
\small
\textit{Reference: Pennington et al. (2014) - GloVe: Global Vectors for Word Representation}
\end{frame}

% ============================================================
% FASTTEXT
% ============================================================
\begin{frame}{FastText: Subword Information 🔤}
\textbf{Word2Vec's limitation: Can't handle unknown words!}

\vspace{1em}

\begin{columns}[T]
\begin{column}{0.48\textwidth}
\textbf{Innovation:}
\begin{itemize}
    \item Represent words as bag of character n-grams
    \item Word = sum of n-gram vectors
    \item Can generate embeddings for OOV words
    \item Captures morphology
\end{itemize}

\vspace{0.5em}
\textbf{Example:}
\begin{itemize}
    \item Word: "where"
    \item 3-grams: <wh, whe, her, ere, re>
    \item $\vec{where} = \sum \vec{n\text{-}gram}$
\end{itemize}
\end{column}

\begin{column}{0.48\textwidth}
\begin{lstlisting}[language=Python]
from gensim.models import FastText

# Train FastText model
model = FastText(
    sentences,
    vector_size=100,
    window=5,
    min_count=1,
    min_n=3,  # min n-gram length
    max_n=6   # max n-gram length
)

# Works for unknown words!
vec = model.wv['unknownword']

# Great for morphologically rich languages
# e.g., German, Finnish, Turkish
\end{lstlisting}

\vspace{0.5em}
\textbf{Advantages:}
\begin{itemize}
    \item Handles rare words better
    \item Morphological awareness
    \item More robust
\end{itemize}
\end{column}
\end{columns}

\vspace{0.5em}
\small
\textit{Reference: Bojanowski et al. (2017) - Enriching Word Vectors with Subword Information}
\end{frame}

% ============================================================
% COMPARISON TABLE
% ============================================================
\begin{frame}{Embedding Methods Comparison 📊}
\begin{center}
\small
\begin{tabular}{lcccccc}
\toprule
\textbf{Method} & \textbf{Year} & \textbf{Type} & \textbf{OOV} & \textbf{Context} & \textbf{Speed} & \textbf{Dims} \\
\midrule
LSA & 1990 & Count & ✗ & Static & Medium & 100-300 \\
LDA & 2003 & Probabilistic & ✗ & Static & Slow & 10-100 \\
Word2Vec & 2013 & Neural & ✗ & Static & Fast & 100-300 \\
GloVe & 2014 & Hybrid & ✗ & Static & Fast & 50-300 \\
FastText & 2017 & Neural & ✓ & Static & Fast & 100-300 \\
\alert{BERT} & 2018 & Transformer & ✓ & \alert{Contextual} & Slow & 768+ \\
\bottomrule
\end{tabular}
\end{center}

\vspace{1em}

\begin{block}{Key Takeaway}
\begin{itemize}
    \item \textbf{Static embeddings:} One vector per word (Word2Vec, GloVe, FastText)
    \item \textbf{Contextual embeddings:} Different vectors based on context (BERT, GPT)
    \item We'll cover contextual embeddings in detail next week!
\end{itemize}
\end{block}
\end{frame}

% ============================================================
% VISUALIZATION EXAMPLE
% ============================================================
\begin{frame}{Visualizing Embeddings with UMAP 🎨}
\begin{columns}[T]
\begin{column}{0.48\textwidth}
\textbf{Example: Word2Vec + UMAP}

\vspace{0.5em}
\begin{lstlisting}[language=Python]
import umap
import matplotlib.pyplot as plt

# Get word vectors
words = ['king', 'queen', 'man',
         'woman', 'cat', 'dog']
vectors = [model[w] for w in words]

# Reduce to 2D
reducer = umap.UMAP(n_components=2)
embedding_2d = reducer.fit_transform(
    vectors
)

# Plot
plt.figure(figsize=(10, 8))
plt.scatter(
    embedding_2d[:, 0],
    embedding_2d[:, 1]
)
for i, word in enumerate(words):
    plt.annotate(
        word,
        (embedding_2d[i, 0],
         embedding_2d[i, 1])
    )
plt.show()
\end{lstlisting}
\end{column}

\begin{column}{0.48\textwidth}
\textbf{What you'll see:}
\begin{itemize}
    \item Semantic clusters
    \item Gender relationships
    \item Animal groupings
    \item Analogical structures
\end{itemize}

\vspace{1em}
\begin{center}
\begin{tikzpicture}[scale=0.9]
    % Axes
    \draw[->] (0,0) -- (5,0) node[right] {UMAP-1};
    \draw[->] (0,0) -- (0,4) node[above] {UMAP-2};

    % Royalty cluster
    \node[circle,fill=purple,inner sep=3pt,label=right:{king}] at (3.5,3) {};
    \node[circle,fill=purple,inner sep=3pt,label=right:{queen}] at (3.8,2.5) {};

    % People cluster
    \node[circle,fill=blue,inner sep=3pt,label=left:{man}] at (2,2.8) {};
    \node[circle,fill=blue,inner sep=3pt,label=left:{woman}] at (2.3,2.3) {};

    % Animal cluster
    \node[circle,fill=green,inner sep=3pt,label=below:{cat}] at (1.5,1) {};
    \node[circle,fill=green,inner sep=3pt,label=below:{dog}] at (2,0.8) {};
\end{tikzpicture}
\end{center}

\vspace{0.5em}
\textbf{Try it yourself in Assignment 3!}
\end{column}
\end{columns}
\end{frame}

% ============================================================
% EVALUATION
% ============================================================
\begin{frame}{How Do We Evaluate Embeddings? 📏}
\begin{columns}[T]
\begin{column}{0.48\textwidth}
\textbf{Intrinsic Evaluation:}
\begin{itemize}
    \item \textbf{Word Similarity:} Compare to human judgments
    \begin{itemize}
        \item WordSim-353
        \item SimLex-999
    \end{itemize}
    \item \textbf{Analogies:} $a:b :: c:d$
    \begin{itemize}
        \item Google Analogy Dataset
        \item Semantic \& Syntactic
    \end{itemize}
    \item \textbf{Clustering:} Do similar words cluster?
\end{itemize}
\end{column}

\begin{column}{0.48\textwidth}
\textbf{Extrinsic Evaluation:}
\begin{itemize}
    \item Use embeddings in downstream tasks
    \item \textbf{Classification:} Sentiment analysis
    \item \textbf{NER:} Named entity recognition
    \item \textbf{QA:} Question answering
    \item \textbf{MT:} Machine translation
\end{itemize}

\vspace{1em}
\alert{Better embeddings $\rightarrow$ Better task performance}
\end{column}
\end{columns}

\vspace{1em}

\begin{block}{Important Note}
No single "best" embedding for all tasks! Choice depends on:
\begin{itemize}
    \item Domain (general vs. specialized)
    \item Task requirements
    \item Computational constraints
\end{itemize}
\end{block}
\end{frame}

% ============================================================
% LIMITATIONS
% ============================================================
\begin{frame}{Limitations of Static Embeddings ⚠️}
\textbf{The fundamental problem: One vector per word}

\vspace{1em}

\begin{exampleblock}{Example: "bank"}
\begin{enumerate}
    \item "I deposited money at the \textbf{bank}" (financial institution)
    \item "We sat by the river \textbf{bank}" (riverside)
\end{enumerate}
Both get the \alert{same vector}! This is wrong.
\end{exampleblock}

\vspace{1em}

\begin{columns}[T]
\begin{column}{0.48\textwidth}
\textbf{Other Issues:}
\begin{itemize}
    \item No compositional semantics
    \item Can't handle phrases well
    \item Bias in embeddings
    \item Limited to training corpus
    \item No common sense reasoning
\end{itemize}
\end{column}

\begin{column}{0.48\textwidth}
\textbf{Solution:}
\begin{itemize}
    \item \alert{Contextual embeddings!}
    \item Different vector for each occurrence
    \item Transformer models (BERT, GPT)
    \item Coming in Week 5-6...
\end{itemize}
\end{column}
\end{columns}
\end{frame}

% ============================================================
% PREVIEW TRANSFORMERS
% ============================================================
\begin{frame}{Preview: Transformer Embeddings 🔮}
\textbf{The next evolution: Context-aware representations}

\vspace{1em}

\begin{columns}[T]
\begin{column}{0.48\textwidth}
\textbf{Static (Word2Vec, GloVe):}
\begin{lstlisting}[language=Python]
# Same vector every time
vec1 = model['bank']
vec2 = model['bank']
# vec1 == vec2 (always!)
\end{lstlisting}

\vspace{0.5em}
\textbf{Problem:} "bank" in different contexts gets same representation
\end{column}

\begin{column}{0.48\textwidth}
\textbf{Contextual (BERT):}
\begin{lstlisting}[language=Python]
from transformers import AutoModel, AutoTokenizer

model = AutoModel.from_pretrained('bert-base-uncased')
tokenizer = AutoTokenizer.from_pretrained('bert-base-uncased')

sent1 = "river bank"
sent2 = "money bank"

# Different vectors!
vec1 = get_embedding(model, tokenizer, sent1, 'bank')
vec2 = get_embedding(model, tokenizer, sent2, 'bank')
# vec1 != vec2
\end{lstlisting}

\vspace{0.5em}
\textbf{Solution:} Different representation based on context!
\end{column}
\end{columns}

\vspace{1em}
\small
\textit{Reference: HuggingFace Course - Chapter 2.2 (Behind the pipeline), Chapter 5.6 (Building a custom dataset)}
\end{frame}

% ============================================================
% INTERACTIVE QUESTION
% ============================================================
\begin{frame}{Discussion: How Do Humans Represent Meaning? 🤔}
\begin{center}
\Large
\textbf{Think about the word "coffee"}
\end{center}

\vspace{1em}

\textbf{What comes to mind?}
\begin{itemize}
    \item Visual: brown liquid, mug, beans
    \item Olfactory: aroma, smell
    \item Gustatory: taste, bitter
    \item Tactile: hot, warm cup
    \item Contextual: morning, caffeine, energy
    \item Emotional: comfort, alertness
    \item Cultural: cafe, social
\end{itemize}

\vspace{1em}

\begin{alertblock}{Key Question}
How does this compare to how Word2Vec represents "coffee"?
\begin{itemize}
    \item Word2Vec: "coffee" = [0.23, -0.45, 0.67, ..., 0.12] (300 numbers)
    \item Learned purely from text co-occurrence
    \item Missing grounding in sensory experience!
\end{itemize}
\end{alertblock}

\small
\textit{This is the "symbol grounding problem" in AI}
\end{frame}

% ============================================================
% PRACTICAL TIPS
% ============================================================
\begin{frame}{Practical Tips for Using Embeddings 💡}
\begin{enumerate}
    \item \textbf{Start with pre-trained models}
    \begin{itemize}
        \item Word2Vec: Google News (3B words)
        \item GloVe: Common Crawl (840B tokens)
        \item FastText: Available in 157 languages
    \end{itemize}

    \item \textbf{Fine-tune if you have domain-specific data}
    \begin{itemize}
        \item Medical texts, legal documents, etc.
        \item Can significantly improve performance
    \end{itemize}

    \item \textbf{Choose dimensions wisely}
    \begin{itemize}
        \item More dims = more expressiveness, more data needed
        \item Typical: 100-300 dimensions
    \end{itemize}

    \item \textbf{Normalize vectors}
    \begin{itemize}
        \item Use cosine similarity, not Euclidean distance
        \item $\text{similarity}(u,v) = \frac{u \cdot v}{||u|| \cdot ||v||}$
    \end{itemize}

    \item \textbf{Be aware of biases}
    \begin{itemize}
        \item Embeddings reflect biases in training data
        \item Gender, race, cultural biases can be encoded
    \end{itemize}
\end{enumerate}
\end{frame}

% ============================================================
% ASSIGNMENT 3
% ============================================================
\begin{frame}{Assignment 3: Exploring Text Embeddings ✍️}
\textbf{Hands-on practice with everything we've learned!}

\vspace{1em}

\textbf{Tasks:}
\begin{enumerate}
    \item \textbf{Build a term-document matrix} from a corpus
    \begin{itemize}
        \item Implement TF-IDF weighting
        \item Visualize with PCA
    \end{itemize}

    \item \textbf{Train your own Word2Vec model}
    \begin{itemize}
        \item Compare CBOW vs Skip-gram
        \item Explore word analogies
    \end{itemize}

    \item \textbf{Visualize embeddings with UMAP}
    \begin{itemize}
        \item Create 2D plots
        \item Identify semantic clusters
    \end{itemize}

    \item \textbf{Topic modeling with LDA}
    \begin{itemize}
        \item Extract topics from documents
        \item Interpret topic coherence
    \end{itemize}

    \item \textbf{Compare embedding methods}
    \begin{itemize}
        \item Word2Vec vs GloVe vs FastText
        \item Evaluate on similarity tasks
    \end{itemize}
\end{enumerate}

\vspace{0.5em}
\textbf{Due:} Check course website for deadline
\end{frame}

% ============================================================
% RESOURCES
% ============================================================
\begin{frame}{Resources \& Further Reading 📚}
\textbf{Key Papers:}
\begin{itemize}
    \item Deerwester et al. (1990) - Latent Semantic Analysis
    \item Blei et al. (2003) - Latent Dirichlet Allocation
    \item Mikolov et al. (2013a,b) - Word2Vec
    \item Pennington et al. (2014) - GloVe
    \item Bojanowski et al. (2017) - FastText
    \item McInnes \& Healy (2018) - UMAP
    \item Boleda (2020) - Distributional Semantics and Linguistic Theory
    \item Grand et al. (2022) - Semantic Projection
\end{itemize}

\vspace{0.5em}

\textbf{Tutorials \& Tools:}
\begin{itemize}
    \item HuggingFace Course: Chapters 2.2, 5.6
    \item Gensim documentation: \url{https://radimrehurek.com/gensim/}
    \item UMAP documentation: \url{https://umap-learn.readthedocs.io/}
    \item Word2Vec tutorial: \url{https://rare-technologies.com/word2vec-tutorial/}
\end{itemize}
\end{frame}

% ============================================================
% SUMMARY
% ============================================================
\begin{frame}{Summary 🎯}
\textbf{What we learned:}

\begin{enumerate}
    \item \textbf{Distributional Hypothesis:} Words in similar contexts have similar meanings

    \item \textbf{Dimensionality Reduction:} PCA (linear), UMAP (non-linear), t-SNE (visualization)

    \item \textbf{Classic Methods:}
    \begin{itemize}
        \item LSA: SVD on term-document matrix
        \item LDA: Probabilistic topic modeling
    \end{itemize}

    \item \textbf{Neural Embeddings:}
    \begin{itemize}
        \item Word2Vec: CBOW \& Skip-gram
        \item GloVe: Global co-occurrence statistics
        \item FastText: Subword information
    \end{itemize}

    \item \textbf{Limitations:} Static embeddings can't handle polysemy

    \item \textbf{Next:} Contextual embeddings (BERT, transformers)
\end{enumerate}

\vspace{1em}
\begin{center}
\Large
\textbf{From words to vectors, from vectors to meaning! 🚀}
\end{center}
\end{frame}

% ============================================================
% NEXT WEEK
% ============================================================
\begin{frame}{Next Week: Neural Language Models 🧠}
\textbf{Building towards modern LLMs}

\vspace{1em}

\begin{itemize}
    \item Recurrent Neural Networks (RNNs)
    \item Long Short-Term Memory (LSTMs)
    \item Sequence-to-Sequence models
    \item Attention mechanisms
    \item The path to transformers...
\end{itemize}

\vspace{2em}

\begin{center}
\Large
\textbf{Questions?} 🙋
\end{center}
\end{frame}

\end{document}
